\subsection{PROPERTIES OF NORMS AND TRACES IN AN ALGEBRA}

PROPOSITION 3. Let A be a K-algebra admitting a finite basis. For an element \(a \in \mathbf{A}\) to be invertible, it is necessary and sufficient that its \(\mathbf{N}_{\mathbf{A} / \mathbf{K}}(a)\) be invertible in \(\mathbf{K}\) .

If \(a\) admits an inverse \(a'\) in \(\mathbf{A}\) , then

\[
\mathrm{N} _ {\mathbf {A} / \mathbf {K}} (a) \mathrm{N} _ {\mathbf {A} / \mathbf {K}} (a ^ {\prime}) = \mathrm{N} _ {\mathbf {A} / \mathbf {K}} (a a ^ {\prime}) = \mathrm{N} _ {\mathbf {A} / \mathbf {K}} (1) = 1
\]

by formula (3) of no. 1. Conversely, if \(\mathbf{N}_{\mathbf{A} / \mathbb{K}}(a)\) is invertible, the endomorphism \(h: x \mapsto ax\) is bijective (§ 8, no. 2, Theorem 1). Then there exists \(a' \in A\) such that \(aa' = 1\) ; then \(h(a'a - 1) = aa'a - a = (aa' - 1)a = 0\) , whence \(aa' = 1\) since \(h\) is injective. Hence \(a'\) is the inverse of \(a\) .

PROPOSITION 4. Let A be a K-algebra admitting a finite basis. For all \(a \in \mathbf{A}\) , \(\operatorname{Pc}_{\mathbf{A} / \mathbf{K}}(a; a) = 0\) .

This follows immediately from the Cayley-Hamilton theorem (§ 8, no. 11, Proposition 20).

PROPOSITION 5. Let A be a K-algebra and m a two-sided ideal of A. Suppose that \(A_{0} = A/m\) is a free K-module of finite dimension n, that there exists an integer r \textgreater{} 0 such that \(m^{r} = \{0\}\) and that \(m^{i-1}/m^{i}\) is a free \(A_{0}\) -module of finite dimension \(s_{i}\) for \(1 \leqslant i \leqslant r\) . Let \(s = s_{1} + \cdots + s_{r}\) and for all \(a \in A\) let \(a_{0}\) denote the class of a mod. m. Then A is a free K-module of dimension ns and, for all \(a \in A\) ,

\[
\operatorname{Tr} _ {\mathbf {A}} (a) = s. \operatorname{Tr} _ {\mathbf {A} _ {0}} (a _ {0}), \quad \mathrm{N} _ {\mathbf {A}} (a) = (\mathrm{N} _ {\mathbf {A} _ {0}} (a _ {0})) ^ {s}\tag{23}
\]

\[
\mathrm{Pc} _ {\mathbf {A}} (a; \mathbf {X}) = (\mathrm{Pc} _ {\mathbf {A} _ {0}} (a _ {0}; \mathbf {X})) ^ {s}.
\]

By virtue of II, § 1, no. 13, Proposition 25, \(\mathfrak{m}^{i - 1} / \mathfrak{m}^{i}\) is a free K-module of dimension \(ns_i\) . Hence Proposition 1 of no. 2 can be applied with \(\mathbf{P}_i = \mathfrak{m}^{i - 1} / \mathfrak{m}^{i}\) ;

this shows in the first place that A is a free K-module of dimension \(n(s_{1} + \cdots + s_{r}) = ns\) . Moreover, the hypothesis implies that the A-module \(P_{i}\) is isomorphic to a direct sum of \(s_{i}\) submodules isomorphic to the A-module \(A_{0}\) ; by Proposition 1 of no. 2, therefore \(\mathrm{N}_{\mathrm{P}_{i}}(a) = \mathrm{N}_{\mathrm{A}_{0}}(a)^{s_{i}}\) ; finally therefore

\[
\mathrm{N} _ {\mathbf {A}} (a) = \mathrm{N} _ {\mathbf {A} _ {0}} (a) ^ {s}.
\]

In this formula \(\mathrm{N}_{\mathbf{A}_{0}}(a)\) is defined by considering \(A_{0}\) as a left A-module and it is equal to the determinant of the K-linear mapping \(x \mapsto ax\) of \(A_{0}\) into itself; but, as \(ax = a_{0}x\) for \(x \in A_{0}\) , \(\mathrm{N}_{\mathbf{A}_{0}}(a) = \mathrm{N}_{\mathbf{A}_{0}}(a_{0})\) , which completes the proof of formula (23) for the norm. The two others are shown analogously.

PROPOSITION 6. Let A be a commutative K-algebra admitting a finite basis over K and V an A-module admitting a finite basis over A. Then V admits a finite basis over K and for every A-endomorphism u of V, if \(u_{\mathbf{K}}\) is the mapping u considered as a K-endomorphism of V,

\[
\operatorname{Tr} (u _ {\mathbb {K}}) = \operatorname{Tr} _ {\mathbf {A} / \mathbb {K}} (\operatorname{Tr} (u)), \quad \det (u _ {\mathbb {K}}) = \mathrm{N} _ {\mathbf {A} / \mathbb {K}} (\det (u))\tag{24}
\]

\[
\operatorname{Pc} (u _ {\mathbf {K}}; \mathbf {X}) = \mathrm{N} _ {\mathbf {A} [ \mathbf {X} ] / \mathbf {K} [ \mathbf {X} ]} (\operatorname{Pc} (u; \mathbf {X})).
\]

Let \((a_{i})_{1\leqslant i\leqslant m}\) be a basis of A over K and \((e_{j})_{1\leqslant j\leqslant n}\) a basis of V over A; then \((a_{i}e_{j})\) is a basis of V over K (II, § 1, no. 13, Proposition 25). On the other hand the third of formulae (24) can be deduced from the second applied to the endomorphism \(\mathbf{X} - \bar{\boldsymbol{u}}\) of the A{[}X{]}-module A{[}X{]} \(\otimes_{\mathbf{A}}\mathrm{V}\) (§ 8, no. 10). It will therefore suffice to show the first two formulae in (24). We shall first establish the following lemma:

Lemma 1. Let \(\mathbf{X}_{ij}\) ( \(1 \leqslant i \leqslant n\) , \(1 \leqslant j \leqslant n\) ) be \(n^2\) indeterminates, \(X\) the square matrix \((\mathbf{X}_{ij})\) of order \(n\) and \(\mathbf{D}(\mathbf{X}_{11}, \ldots, \mathbf{X}_{nn}) \in \mathbf{Z}[\mathbf{X}_{11}, \ldots, \mathbf{X}_{nn}]\) the determinant \(\det(X)\) . On the other hand let A be a commutative ring, \(\mathbf{M}_{ij}\) ( \(1 \leqslant i \leqslant n\) , \(1 \leqslant j \leqslant n\) ) \(n^2\) matrices of order \(m\) over A, which are pairwise permutable, and M the square matrix of order \(mn\) over A which can be expressed as a square matrix of matrices (II, § 10, no. 7)

\[
M = \left( \begin{array}{c c c c} M _ {1 1} & M _ {1 2} & \ldots & M _ {1 n} \\ M _ {2 1} & M _ {2 2} & \ldots & M _ {2 n} \\ \cdot & \cdot & \cdot & \cdot \\ M _ {n 1} & M _ {n 2} & \ldots & M _ {n n} \end{array} \right)
\]

Then the determinant of M is equal to the determinant of the square matrix

\[
\mathrm{D} (M _ {1 1}, \dots , M _ {n n})
\]

of order m.

We proceed by induction on \(n\) , the cases \(n = 0\) and \(n = 1\) being trivial.

Let \(Z\) be a new indeterminate and \(N_{ij}\) the matrix \(M_{ij} + \delta_{ij}ZI_m\) ( \(\delta_{ij}\) the Kronecker index). If \(D^{ij}(X_{11},\ldots,X_{nn})\) is the cofactor of \(X_{ij}\) in the matrix \(X\) (§ 8, no. 6), then

\[
\mathrm{X} _ {j l} \mathrm{D} ^ {k l} (\mathrm{X} _ {1 1}, \dots , \mathrm{X} _ {n n}) = \delta_ {j k} \mathrm{D} (\mathrm{X} _ {1 1}, \dots , \mathrm{X} _ {n n})\tag{25}
\]

(§ 8, no. 6, formula (28)). We write \(N_{ij}' = \mathbf{D}^{ij}(N_{11}, \ldots, N_{nn})\) , which is a square matrix of order \(m\) over \(\mathrm{A}[Z]\) and consider the product \(N.U\) , where

\[
U = \left( \begin{array}{c c c c} N _ {1 1} ^ {\prime} & 0 & \ldots & 0 \\ N _ {1 2} ^ {\prime} & I _ {m} & \ldots & 0 \\ \cdot & \cdot & \cdot & \cdot \\ N _ {1 n} ^ {\prime} & 0 & \ldots & I _ {m} \end{array} \right),
\]

\[
N = \left( \begin{array}{c c c c} N _ {1 1} & N _ {1 2} & \ldots & N _ {1 n} \\ N _ {2 1} & N _ {2 2} & \ldots & N _ {2 n} \\ \cdot & \cdot & \cdot & \cdot \\ N _ {n 1} & N _ {n 2} & \ldots & N _ {n n} \end{array} \right)
\]

Performing this product in blocks (II, § 10, no. 5) and using formulae (25), we obtain

\[
N. U = \left( \begin{array}{c c c c} P & N _ {1 2} & \ldots & N _ {1 n} \\ 0 & N _ {2 2} & \ldots & N _ {2 n} \\ \cdot & \cdot & \cdot & \cdot \\ 0 & N _ {n 2} & \ldots & N _ {n n} \end{array} \right)
\]

where we have written \(P = \mathrm{D}(N_{11},\ldots ,N_{nn})\) .Let

\[
Q = \left( \begin{array}{c c c c} N _ {2 2} & \ldots & N _ {2 n} \\ \cdot & \cdot & \cdot & \cdot \\ N _ {n 2} & \ldots & N _ {n n} \end{array} \right)
\]

which is a matrix of order \(m(n - 1)\) ; then (§8, no. 6, formula (31)) (det \(N\) ) ( \(\det U\) ) = ( \(\det P\) ) ( \(\det Q\) ) and \(\det U = \det N_{11}'\) . But by the induction hypothesis, \(\det Q = \det(\mathrm{D}^{11}(N_{11},\ldots ,N_{nn})) = \det N_{11}'\) and by virtue of the definition of the \(N_{ij}\) , clearly \(\det Q\) is a polynomial in A{[}Z{]} of degree \(m(n - 1)\) , whose term in \(Z^{m(n - 1)}\) has coefficient 1; it follows immediately that \(\det Q\) is not a divisor of zero in the graded algebra A{[}Z{]}. We therefore conclude that \(\det N = \det(\mathrm{D}(N_{11},\ldots ,N_{nn}))\) in A{[}Z{]}; if we substitute 0 for Z in these polynomials, then \(\det M = \det(\mathrm{D}(M_{11},\ldots ,M_{nn}))\) .

Having shown this lemma, the K-module V is the direct sum of the K-modules \(Ae_{j}\) ( \(1 \leqslant j \leqslant n\) ); we write \(u(e_{j}) = \sum_{k=1}^{n} c_{jk} e_{k}\) . For every element \(xe_{j} \in Ae_{j}\) , with \(x \in A\) , the component of \(u(xe_{j})\) in \(Ae_{k}\) is \(xc_{jk} e_{k}\) ; it follows that the matrix of \(u_{k}\) with respect to the basis \((a_{i} e_{j})\) of the K-module V can be expressed in the form of a square matrix of matrices \((M_{jk})\) , where \(M_{jk}\) is the matrix of the K-linear mapping \(xe_{j} \mapsto xc_{jk} e_{k}\) of \(Ae_{j}\) into \(Ae_{k}\) with respect to the bases \((a_{i} e_{j})_{1 \leqslant i \leqslant m}\) and \((a_{i} e_{k})_{1 \leqslant i \leqslant m}\) of these two K-modules (II, § 10, no. 5). If for all \(t \in A\) , \(M(t)\) denotes the matrix, with respect to the basis \((a_{i})_{1 \leqslant i \leqslant m}\) of A over K, of the endomorphism \(x \mapsto xt\) of A, then \(M_{jk} = M(c_{jk})\) ; as \(t \mapsto M(t)\) is a ring homomorphism the matrices \(M_{jk}\) are permutable with one another. Then

\[
\det u _ {\mathrm{K}} = \det (\mathrm{D} (M _ {1 1}, \dots , M _ {n n}))
\]

by Lemma 1. But as \(t \mapsto M(t)\) is a ring homomorphism, \(\mathrm{D}(M_{11},\ldots ,M_{nn})\) is the matrix of the K-endomorphism \(x \mapsto x.\det (c_{jk})\) of A with respect to basis \((a_t)\) ; by definition its determinant is therefore \(\mathrm{N}_{\mathbf{A} / \mathbf{K}}(\det (u))\) , which proves the second formula of (24). On the other hand,

\[
\operatorname{Tr} \left(u _ {\mathrm{K}}\right) = \sum_ {j = 1} ^ {n} \operatorname{Tr} \left(M _ {j j}\right) = \sum_ {j = 1} ^ {n} \operatorname{Tr} _ {\mathrm{A} / \mathrm{K}} \left(c _ {j j}\right) = \operatorname{Tr} _ {\mathrm{A} / \mathrm{K}} \left(\sum_ {j = 1} ^ {n} c _ {j j}\right) = \operatorname{Tr} _ {\mathrm{A} / \mathrm{K}} (\operatorname{Tr} (u))
\]

and the proof of Proposition 6 is complete.

COROLLARY. Let A be a commutative K-algebra admitting a finite basis over K and B an A-algebra admitting a finite basis over A. Then B admits a finite basis over K, and, for all \(b \in B\) (``transitivity formulae'')

\[
\operatorname{Tr} _ {\mathbf {B} / \mathbf {K}} (b) = \operatorname{Tr} _ {\mathbf {A} / \mathbf {K}} (\operatorname{Tr} _ {\mathbf {B} / \mathbf {A}} (b)), \quad \mathrm{N} _ {\mathbf {B} / \mathbf {K}} (b) = \mathrm{N} _ {\mathbf {A} / \mathbf {K}} (\mathrm{N} _ {\mathbf {B} / \mathbf {A}} (b))\tag{26}
\]

\[
\mathrm{Pc} _ {\mathbf {B} / \mathbf {K}} (b; \mathbf {X}) = \mathrm{N} _ {\mathbf {A} [ \mathbf {X} ] / \mathbf {K} [ \mathbf {X} ]} (\mathrm{Pc} _ {\mathbf {B} / \mathbf {A}} (b; \mathbf {X})).
\]

This follows immediately from Proposition 6, setting \(\mathbf{V} = \mathbf{B}\) and \(u(x) = bx\) .

Remark. Suppose that the homomorphism \(\lambda \mapsto \lambda, 1\) of K into A is injective and let K be identified with its image in A; suppose that A admits a finite basis \((e_i)_{1 \leqslant i \leqslant n}\) as a K-module. Let \(s\) be an automorphism of A such that \(s(\mathbf{K}) = \mathbf{K}\) . Let \(a\) be an element of A; then, by transporting the structure

\begin{enumerate}
\def\labelenumi{(\arabic{enumi})}
\setcounter{enumi}{26}
\tightlist
\item
\end{enumerate}

\[
\operatorname{Tr} _ {\mathbf {A} / \mathbf {K}} (s (a)) = s (\operatorname{Tr} _ {\mathbf {A} / \mathbf {K}} (a))\tag{28}
\]

\[
\mathrm{N} _ {\mathbf {A} / \mathbf {K}} (s (a)) = s (\mathrm{N} _ {\mathbf {A} / \mathbf {K}} (a)).
\]

*Consider also a derivation D of A (§ 10, no. 2) such that D(K) ⊂ K and write \(D (e _ {i}) = \sum_ {j = 1} ^ {n} e _ {j} \mu_ {j i}\) where \(\mu_ {j i} \in \mathbf {K}\); write

\[
a e _ {i} = \sum_ {j = 1} ^ {n} e _ {j} \lambda_ {j i} \quad \text { with } \quad \lambda_ {j i} \in \mathbf {K}.
\]

Then

\[
\mathrm{D} (a) e _ {i} + a \mathrm{D} (e _ {i}) = \mathrm{D} (a e _ {i}) = \sum_ {j = 1} ^ {n} (\mathrm{D} (e _ {j}) \lambda_ {j i} + e _ {j} \mathrm{D} (\lambda_ {j i})).
\]

It follows that

\[
\mathbf {D} (a) e _ {i} = \sum_ {j = 1} ^ {n} e _ {j} \nu_ {j i}
\]

with \(\nu_{ji} = \mathrm{D}(\lambda_{ji}) + \sum_{k=1}^{n} (\mu_{jk}\lambda_{ki} - \lambda_{jk}\mu_{ki})\) . As \(\sum_{i,k} (\mu_{ik}\lambda_{ki} - \lambda_{ik}\mu_{ki}) = 0\) , there-

fore \(\mathrm{Tr}_{\mathbf{A} / \mathbb{K}}(\mathbf{D}(a)) = \sum_{i = 1}^{n}\mathbf{D}(\lambda_{ii})\) , in other words

\[
\operatorname{Tr} _ {\mathbf {A} / \mathbf {K}} (\mathbf {D} (a)) = \mathbf {D} (\operatorname{Tr} _ {\mathbf {A} / \mathbf {K}} (a)). *\tag{29}
\]

